\documentclass[12pt,a4paper]{article}

\usepackage[utf8]{inputenc}
\usepackage[T1]{fontenc}
\usepackage[french]{babel}
\usepackage[margin=2.5cm]{geometry}
\usepackage{listings}
\usepackage{xcolor}
\usepackage{hyperref}
\usepackage{parskip}
\usepackage{booktabs}
\usepackage{fancyhdr}
\usepackage{lmodern}

\setlength{\headheight}{14.5pt}

% ── style des blocs de code ───────────────────────────────────────────────────
\definecolor{codebg}{HTML}{F5F5F5}
\definecolor{keyword}{HTML}{0000CC}
\definecolor{comment}{HTML}{007700}
\definecolor{string}{HTML}{CC0000}
\definecolor{linenum}{HTML}{888888}

\lstdefinestyle{cstyle}{
    language=C,
    backgroundcolor=\color{codebg},
    basicstyle=\ttfamily\small,
    keywordstyle=\color{keyword}\bfseries,
    commentstyle=\color{comment}\itshape,
    stringstyle=\color{string},
    numberstyle=\tiny\color{linenum},
    numbers=left,
    numbersep=8pt,
    stepnumber=1,
    breaklines=true,
    breakatwhitespace=false,
    tabsize=4,
    showstringspaces=false,
    frame=single,
    framesep=4pt,
    xleftmargin=16pt,
    xrightmargin=4pt,
    captionpos=b,
}

\lstset{style=cstyle}

% ── en-tête et pied de page ───────────────────────────────────────────────────
\pagestyle{fancy}
\fancyhf{}
\rhead{Sockets 3, Programmation Système}
\lhead{Chat Hub Multi-clients}
\cfoot{\thepage}

% ── document ──────────────────────────────────────────────────────────────────
\begin{document}

\begin{titlepage}
    \centering
    \vspace*{2cm}
    {\Huge\bfseries Chat Hub Multi-clients\par}
    \vspace{0.5cm}
    {\Large Programmation Système en C, Sockets 3\par}
    \vspace{2cm}
    \rule{\linewidth}{0.4pt}\\[0.4cm]
    {\large Serveur de chat multi-clients avec protocole texte et threads détachés\par}
    \rule{\linewidth}{0.4pt}
    \vfill
    {\large Ilyass Bougati\par}
    \vspace{0.3cm}
    {\large \today\par}
\end{titlepage}

\tableofcontents
\newpage

% ── Introduction ──────────────────────────────────────────────────────────────
\section{Introduction}

Ce projet est l'évolution naturelle du chat full-duplex (\emph{Sockets 2}).
Les projets précédents ne mettaient en relation que \emph{deux} interlocuteurs
à la fois. Ici, le serveur devient un \textbf{hub} : il accepte jusqu'à
\texttt{MAX\_CLIENTS} (64) connexions simultanées et redistribue les messages
à l'ensemble des participants connectés.

Deux nouveautés architecturales accompagnent ce passage au multi-clients :

\begin{itemize}
    \item Un \textbf{protocole texte structuré} (verbes \texttt{NICK},
          \texttt{MSG}, \texttt{PRIV}, \texttt{LIST}, \texttt{QUIT}\ldots)
          remplace le flux brut des projets précédents.
    \item Un \textbf{fil d'exécution détaché} est créé pour chaque client,
          accompagné d'un tableau partagé protégé par un mutex.
\end{itemize}

Le code est organisé en trois fichiers : \texttt{common.h} (protocole partagé
et utilitaires réseau), \texttt{server.c} (hub) et \texttt{client.c}
(interface utilisateur).

% ── Protocole applicatif ──────────────────────────────────────────────────────
\section{Protocole applicatif}

\subsection{Format des messages}

Tous les messages sont des \textbf{lignes de texte} terminées par
\texttt{\textbackslash n}. Le premier mot de chaque ligne est le
\emph{verbe} ; les arguments suivent, séparés par des espaces. Le lecteur de
lignes accepte aussi le terminateur CRLF (\texttt{\textbackslash r\textbackslash n})
pour la tolérance avec des outils comme \texttt{telnet} ou \texttt{netcat}.

\subsection{Verbes client $\to$ serveur}

\begin{center}
\begin{tabular}{@{}lll@{}}
\toprule
Verbe & Syntaxe & Description \\
\midrule
\texttt{NICK} & \texttt{NICK <pseudo>} & Choisir un pseudonyme (handshake) \\
\texttt{MSG}  & \texttt{MSG <texte>}   & Diffuser un message à tous \\
\texttt{PRIV} & \texttt{PRIV <pseudo> <texte>} & Message privé \\
\texttt{LIST} & \texttt{LIST}          & Demander la liste des connectés \\
\texttt{QUIT} & \texttt{QUIT}          & Quitter proprement \\
\bottomrule
\end{tabular}
\end{center}

\subsection{Verbes serveur $\to$ client}

\begin{center}
\begin{tabular}{@{}lll@{}}
\toprule
Verbe & Syntaxe & Description \\
\midrule
\texttt{OK}    & \texttt{OK}                     & Handshake accepté \\
\texttt{ERR}   & \texttt{ERR <raison>}            & Erreur applicative \\
\texttt{MSG}   & \texttt{MSG <from> <texte>}      & Message reçu de \texttt{from} \\
\texttt{PRIV}  & \texttt{PRIV <from> <texte>}     & Message privé reçu \\
\texttt{JOIN}  & \texttt{JOIN <pseudo>}            & Notification d'arrivée \\
\texttt{LEAVE} & \texttt{LEAVE <pseudo>}           & Notification de départ \\
\texttt{LIST}  & \texttt{LIST <p1> <p2> \ldots}   & Liste des connectés \\
\texttt{SYS}   & \texttt{SYS <texte>}             & Message système \\
\bottomrule
\end{tabular}
\end{center}

\subsection{Handshake de connexion}

À la connexion, le client doit s'identifier avant de pouvoir envoyer quoi que
ce soit :

\begin{enumerate}
    \item Le client envoie \texttt{NICK <pseudo>}.
    \item Le serveur répond \texttt{OK} si le pseudo est valide et disponible,
          ou \texttt{ERR <raison>} sinon.
    \item Le client peut réessayer avec un autre pseudo en cas d'erreur.
\end{enumerate}

Le pseudo n'accepte que des caractères alphanumériques, \texttt{\_} et
\texttt{-} (au plus 31 caractères). Cette validation est effectuée par la
fonction \texttt{valid\_name()} côté serveur.

% ── Architecture ──────────────────────────────────────────────────────────────
\section{Architecture du serveur}

\subsection{Modèle de threading}

Le serveur suit le modèle classique \emph{un fil par client} :

\begin{center}
\begin{tabular}{@{}lll@{}}
\toprule
Fil & Rôle & Cycle de vie \\
\midrule
Principal & \texttt{accept()} en boucle & Durée de vie du processus \\
\texttt{client\_thread} & Gérer un client (lecture + dispatch) & Durée de la connexion \\
\bottomrule
\end{tabular}
\end{center}

Les fils clients sont créés avec \texttt{pthread\_create} puis
\textbf{détachés} immédiatement via \texttt{pthread\_detach}. Le serveur
n'attend donc jamais leur terminaison (\texttt{pthread\_join}) : chaque fil
libère ses ressources automatiquement en sortant.

\subsection{Table des clients et exclusion mutuelle}

L'état partagé est un tableau global \texttt{clients[MAX\_CLIENTS]} de type
\texttt{client\_t} :

\begin{lstlisting}[caption={Structure d'un client enregistré}]
typedef struct {
    int    fd;
    char   name[NAME_SIZE];
    time_t joined_at;
    int    in_use;
} client_t;
\end{lstlisting}

Toute lecture ou écriture de ce tableau est protégée par
\texttt{clients\_mu} (un \texttt{pthread\_mutex\_t} statique). Le journal
d'événements dispose de son propre mutex \texttt{log\_mu} pour éviter
l'entrelacement des lignes entre fils.

\subsection{Diffusion sans blocage}

La fonction \texttt{broadcast()} est conçue pour qu'un client lent ne bloque
pas les autres :

\begin{enumerate}
    \item Elle prend un \emph{instantané} des descripteurs actifs sous verrou.
    \item Elle relâche le verrou avant d'appeler \texttt{send\_line} sur chaque
          descripteur.
\end{enumerate}

De plus, chaque descripteur de client se voit appliquer une option
\texttt{SO\_SNDTIMEO} de 5 secondes dès l'acceptation de la connexion.
Un envoi qui dépasse ce délai échoue immédiatement plutôt que de bloquer
indéfiniment le fil de diffusion.

\subsection{Lecteur de lignes bufférisé (\texttt{line\_reader\_t})}

TCP est un protocole de flux : un seul \texttt{recv} peut ramener plusieurs
messages, ou un message peut arriver en plusieurs fragments. La structure
\texttt{line\_reader\_t} définie dans \texttt{common.h} gère cette
problématique avec un tampon interne de \texttt{2\,×\,BUF\_SIZE} octets :

\begin{lstlisting}[caption={\texttt{lr\_next}, lecture d'une ligne complète}]
static inline int lr_next(line_reader_t *lr, char *out, size_t out_cap) {
    for (;;) {
        char *nl = memchr(lr->buf, '\n', lr->len);
        if (nl) {
            /* Une ligne complete est disponible : la copier et decaler. */
            ...
            return 1;
        }
        if (lr->len >= sizeof(lr->buf))
            lr->len = 0; /* ligne trop longue : on la jette */
        ssize_t n = recv(lr->fd, lr->buf + lr->len, ...);
        if (n <= 0) return n; /* fermeture ou erreur */
        lr->len += (size_t)n;
    }
}
\end{lstlisting}

\texttt{lr\_next} retourne 1 si une ligne a été produite, 0 en cas de
fermeture propre par le pair, et -1 en cas d'erreur.

% ── Architecture du client ────────────────────────────────────────────────────
\section{Architecture du client}

Le client reprend le modèle à deux fils introduit dans \emph{Sockets 2} :

\begin{itemize}
    \item Le \textbf{fil principal} lit l'entrée standard et envoie des
          commandes au serveur.
    \item Le \textbf{fil de réception} (\texttt{recv\_thread}) lit les lignes
          entrantes et les affiche, en effaçant d'abord le prompt avec la
          séquence ANSI \texttt{\textbackslash r\textbackslash 033[K}.
\end{itemize}

La résolution du nom d'hôte utilise \texttt{getaddrinfo} (plutôt que
\texttt{inet\_pton}), ce qui permet de se connecter par nom DNS en plus de
l'adresse IP.

\subsection{Commandes utilisateur}

Le client traduit les commandes saisies par l'utilisateur en verbes de
protocole :

\begin{center}
\begin{tabular}{@{}ll@{}}
\toprule
Saisie utilisateur & Verbe envoyé \\
\midrule
\texttt{/users} & \texttt{LIST} \\
\texttt{/msg <pseudo> <texte>} & \texttt{PRIV <pseudo> <texte>} \\
\texttt{/quit} & \texttt{QUIT} \\
\textit{(texte libre)} & \texttt{MSG <texte>} \\
\bottomrule
\end{tabular}
\end{center}

% ── Appels système et API ─────────────────────────────────────────────────────
\section{Appels système et API utilisés}

\begin{description}
    \item[\texttt{socket} / \texttt{bind} / \texttt{listen} / \texttt{accept} / \texttt{connect}]
        Infrastructure TCP standard, identique aux projets précédents.
    \item[\texttt{send} / \texttt{recv}]
        Envoi et réception de données. \texttt{send\_all} boucle sur
        \texttt{send} pour garantir l'envoi complet de chaque ligne.
        \texttt{MSG\_NOSIGNAL} empêche la génération d'un \texttt{SIGPIPE}
        si le pair a fermé la connexion.
    \item[\texttt{pthread\_create} / \texttt{pthread\_detach}]
        Un fil par client, détaché pour que sa terminaison libère ses
        ressources sans \texttt{pthread\_join}.
    \item[\texttt{pthread\_mutex\_lock} / \texttt{pthread\_mutex\_unlock}]
        Protection du tableau \texttt{clients[]} et du fichier de journal.
    \item[\texttt{setsockopt(SO\_REUSEADDR)}]
        Redémarrage rapide du serveur sans attendre l'expiration de
        \texttt{TIME\_WAIT}.
    \item[\texttt{setsockopt(SO\_SNDTIMEO)}]
        Délai maximal de 5 secondes par envoi, pour éviter qu'un client
        lent ne bloque les diffusions vers les autres.
    \item[\texttt{shutdown(SHUT\_RDWR)}]
        Déconnexion propre : débloque les \texttt{recv} bloquants des deux
        côtés avant la fermeture du descripteur.
    \item[\texttt{signal(SIGPIPE, SIG\_IGN)}]
        Ignorer \texttt{SIGPIPE} pour que l'écriture sur une socket fermée
        retourne une erreur plutôt que de tuer le processus.
    \item[\texttt{signal(SIGINT, ...)}]
        Fermeture propre du serveur (descripteur d'écoute + fichier log)
        sur Ctrl+C.
    \item[\texttt{getaddrinfo} / \texttt{freeaddrinfo}]
        Résolution DNS côté client, permettant de spécifier un nom d'hôte
        en plus d'une adresse IP.
    \item[\texttt{localtime\_r} / \texttt{strftime}]
        Horodatage thread-safe des entrées du journal (version réentrante
        de \texttt{localtime}).
\end{description}

% ── Utilisation ───────────────────────────────────────────────────────────────
\section{Utilisation}

\begin{lstlisting}[language=bash, style=cstyle]
# Terminal 1 : lancer le hub sur le port 5555 avec journal
./chat_serverd --port 5555 --log chat.log

# Terminal 2, 3, ... : connecter des clients
./chat_client --server 127.0.0.1 --port 5555
\end{lstlisting}

Une fois connecté, les commandes disponibles sont :

\begin{lstlisting}[language=bash, style=cstyle]
/users               # lister les utilisateurs connectes
/msg alice bonjour   # envoyer un message prive a alice
/quit                # quitter proprement
bonjour a tous       # diffuser a tout le salon
\end{lstlisting}

La compilation requiert l'option \texttt{-pthread} (déjà présente dans le
\texttt{Makefile}) pour lier la bibliothèque POSIX Threads.

% ── Code source ───────────────────────────────────────────────────────────────
\section{Code source}

\subsection{\texttt{common.h}}

\begin{lstlisting}[caption={\texttt{common.h}, protocole, constantes et utilitaires réseau}]
#ifndef HUB_COMMON_H
#define HUB_COMMON_H

#include <stddef.h>
#include <string.h>
#include <sys/socket.h>

#define BUF_SIZE     1024
#define NAME_SIZE      32
#define MAX_CLIENTS    64

/* Wire protocol -- line-based, '\n' terminated, first token is the verb.
 *
 * Client -> Server:
 *   NICK <name>
 *   MSG  <text>
 *   PRIV <name> <text>
 *   LIST
 *   QUIT
 *
 * Server -> Client:
 *   OK
 *   ERR  <reason>
 *   MSG  <from> <text>
 *   PRIV <from> <text>
 *   JOIN <name>
 *   LEAVE <name>
 *   LIST <name1> <name2> ...
 *   SYS  <text>
 */

typedef struct {
    int    fd;
    char   buf[BUF_SIZE * 2];
    size_t len;
} line_reader_t;

static inline void lr_init(line_reader_t *lr, int fd) {
    lr->fd = fd;
    lr->len = 0;
}

static inline int lr_next(line_reader_t *lr, char *out, size_t out_cap) {
    for (;;) {
        char *nl = memchr(lr->buf, '\n', lr->len);
        if (nl) {
            size_t line_len = (size_t)(nl - lr->buf);
            size_t copy = line_len < out_cap - 1 ? line_len : out_cap - 1;
            memcpy(out, lr->buf, copy);
            out[copy] = '\0';
            if (copy > 0 && out[copy - 1] == '\r')
                out[copy - 1] = '\0';
            size_t consumed = line_len + 1;
            memmove(lr->buf, lr->buf + consumed, lr->len - consumed);
            lr->len -= consumed;
            return 1;
        }
        if (lr->len >= sizeof(lr->buf))
            lr->len = 0;
        ssize_t n = recv(lr->fd, lr->buf + lr->len,
                         sizeof(lr->buf) - lr->len, 0);
        if (n == 0) return 0;
        if (n  < 0) return -1;
        lr->len += (size_t)n;
    }
}

static inline int send_all(int fd, const char *data, size_t len) {
    while (len > 0) {
        ssize_t n = send(fd, data, len, MSG_NOSIGNAL);
        if (n <= 0) return -1;
        data += n;
        len  -= (size_t)n;
    }
    return 0;
}

static inline int send_line(int fd, const char *line) {
    size_t len = strlen(line);
    if (send_all(fd, line, len) < 0) return -1;
    return send_all(fd, "\n", 1);
}

#endif
\end{lstlisting}

\subsection{\texttt{server.c}}

\begin{lstlisting}[caption={\texttt{server.c}, hub de chat multi-clients}]
#include <stdio.h>
#include <stdlib.h>
#include <stdint.h>
#include <stdarg.h>
#include <string.h>
#include <unistd.h>
#include <signal.h>
#include <pthread.h>
#include <arpa/inet.h>
#include <time.h>
#include <ctype.h>
#include <errno.h>

#include "common.h"

typedef struct {
    int    fd;
    char   name[NAME_SIZE];
    time_t joined_at;
    int    in_use;
} client_t;

static client_t        clients[MAX_CLIENTS];
static pthread_mutex_t clients_mu = PTHREAD_MUTEX_INITIALIZER;
static int             server_fd  = -1;
static FILE           *log_fp     = NULL;
static pthread_mutex_t log_mu     = PTHREAD_MUTEX_INITIALIZER;

static void log_message(const char *fmt, ...) {
    if (!log_fp) return;
    time_t t = time(NULL);
    struct tm tm;
    localtime_r(&t, &tm);
    char ts[32];
    strftime(ts, sizeof(ts), "%Y-%m-%d %H:%M:%S", &tm);
    pthread_mutex_lock(&log_mu);
    fprintf(log_fp, "%s ", ts);
    va_list ap;
    va_start(ap, fmt);
    vfprintf(log_fp, fmt, ap);
    va_end(ap);
    fputc('\n', log_fp);
    fflush(log_fp);
    pthread_mutex_unlock(&log_mu);
}

static int valid_name(const char *s) {
    if (!*s) return 0;
    if (strlen(s) >= NAME_SIZE) return 0;
    for (const char *p = s; *p; p++) {
        if (!(isalnum((unsigned char)*p) || *p == '_' || *p == '-'))
            return 0;
    }
    return 1;
}

static int register_client(int fd, const char *name) {
    pthread_mutex_lock(&clients_mu);
    int slot = -1;
    for (int i = 0; i < MAX_CLIENTS; i++) {
        if (clients[i].in_use && strcmp(clients[i].name, name) == 0) {
            pthread_mutex_unlock(&clients_mu);
            return -1;
        }
        if (slot < 0 && !clients[i].in_use)
            slot = i;
    }
    if (slot < 0) {
        pthread_mutex_unlock(&clients_mu);
        return -2;
    }
    clients[slot].fd        = fd;
    clients[slot].joined_at = time(NULL);
    clients[slot].in_use    = 1;
    strncpy(clients[slot].name, name, NAME_SIZE - 1);
    clients[slot].name[NAME_SIZE - 1] = '\0';
    pthread_mutex_unlock(&clients_mu);
    return slot;
}

static void unregister_client(int slot) {
    pthread_mutex_lock(&clients_mu);
    clients[slot].in_use  = 0;
    clients[slot].fd      = -1;
    clients[slot].name[0] = '\0';
    pthread_mutex_unlock(&clients_mu);
}

static void broadcast(const char *line, int except_fd) {
    int fds[MAX_CLIENTS];
    int n = 0;
    pthread_mutex_lock(&clients_mu);
    for (int i = 0; i < MAX_CLIENTS; i++) {
        if (clients[i].in_use && clients[i].fd != except_fd)
            fds[n++] = clients[i].fd;
    }
    pthread_mutex_unlock(&clients_mu);
    for (int i = 0; i < n; i++)
        (void)send_line(fds[i], line);
}

static int send_private(const char *target, const char *line) {
    int fd = -1;
    pthread_mutex_lock(&clients_mu);
    for (int i = 0; i < MAX_CLIENTS; i++) {
        if (clients[i].in_use && strcmp(clients[i].name, target) == 0) {
            fd = clients[i].fd;
            break;
        }
    }
    pthread_mutex_unlock(&clients_mu);
    if (fd < 0) return -1;
    return send_line(fd, line);
}

static void handle_list(int fd) {
    char reply[BUF_SIZE];
    size_t off = 0;
    off += (size_t)snprintf(reply + off, sizeof(reply) - off, "LIST");
    pthread_mutex_lock(&clients_mu);
    for (int i = 0; i < MAX_CLIENTS && off < sizeof(reply) - 1; i++) {
        if (clients[i].in_use) {
            int w = snprintf(reply + off, sizeof(reply) - off,
                             " %s", clients[i].name);
            if (w < 0) break;
            off += (size_t)w;
        }
    }
    pthread_mutex_unlock(&clients_mu);
    send_line(fd, reply);
}

static void *client_thread(void *arg) {
    int fd = (int)(intptr_t)arg;
    line_reader_t lr;
    lr_init(&lr, fd);

    char line[BUF_SIZE];
    char name[NAME_SIZE] = {0};
    int  slot = -1;

    while (slot < 0) {
        int r = lr_next(&lr, line, sizeof(line));
        if (r <= 0) goto done;
        if (strncmp(line, "NICK ", 5) != 0) {
            send_line(fd, "ERR expected NICK first");
            continue;
        }
        const char *requested = line + 5;
        if (!valid_name(requested)) {
            send_line(fd, "ERR invalid name (alnum, _-, max 31 chars)");
            continue;
        }
        int s = register_client(fd, requested);
        if (s == -1) { send_line(fd, "ERR name already taken"); continue; }
        if (s == -2) { send_line(fd, "ERR server full");        goto done; }
        slot = s;
        strncpy(name, requested, NAME_SIZE - 1);
        name[NAME_SIZE - 1] = '\0';
        send_line(fd, "OK");
    }

    char ann[BUF_SIZE];
    snprintf(ann, sizeof(ann), "JOIN %s", name);
    broadcast(ann, fd);
    log_message("[JOIN] %s", name);

    for (;;) {
        int r = lr_next(&lr, line, sizeof(line));
        if (r <= 0) break;
        if (strcmp(line, "QUIT") == 0) break;
        if (strcmp(line, "LIST") == 0) { handle_list(fd); continue; }
        if (strncmp(line, "MSG ", 4) == 0) {
            char out[BUF_SIZE];
            snprintf(out, sizeof(out), "MSG %s %s", name, line + 4);
            broadcast(out, fd);
            log_message("[MSG] %s: %s", name, line + 4);
            continue;
        }
        if (strncmp(line, "PRIV ", 5) == 0) {
            const char *rest = line + 5;
            const char *sp   = strchr(rest, ' ');
            if (!sp || sp == rest) {
                send_line(fd, "ERR usage: PRIV <name> <text>");
                continue;
            }
            size_t tlen = (size_t)(sp - rest);
            if (tlen >= NAME_SIZE) {
                send_line(fd, "ERR target name too long");
                continue;
            }
            char target[NAME_SIZE];
            memcpy(target, rest, tlen);
            target[tlen] = '\0';
            if (strcmp(target, name) == 0) {
                send_line(fd, "ERR cannot PRIV yourself");
                continue;
            }
            char out[BUF_SIZE];
            snprintf(out, sizeof(out), "PRIV %s %s", name, sp + 1);
            if (send_private(target, out) < 0) {
                char err[BUF_SIZE];
                snprintf(err, sizeof(err), "ERR no such user: %s", target);
                send_line(fd, err);
            } else {
                log_message("[PRIV] %s -> %s: %s", name, target, sp + 1);
            }
            continue;
        }
        send_line(fd, "ERR unknown command");
    }

done:
    if (slot >= 0) {
        unregister_client(slot);
        char leave[BUF_SIZE];
        snprintf(leave, sizeof(leave), "LEAVE %s", name);
        broadcast(leave, fd);
        log_message("[LEAVE] %s", name);
    }
    shutdown(fd, SHUT_RDWR);
    close(fd);
    return NULL;
}

static void handle_sigint(int sig) {
    (void)sig;
    printf("\nShutting down Chat Hub...\n");
    if (server_fd != -1) close(server_fd);
    if (log_fp) fclose(log_fp);
    _exit(0);
}

int main(int argc, char *argv[]) {
    int         port     = 5555;
    const char *log_path = NULL;

    for (int i = 1; i < argc; i++) {
        if (strcmp(argv[i], "--port") == 0 && i + 1 < argc)
            port = atoi(argv[++i]);
        else if (strcmp(argv[i], "--log") == 0 && i + 1 < argc)
            log_path = argv[++i];
        else {
            fprintf(stderr, "Usage: %s [--port <port>] [--log <file>]\n",
                    argv[0]);
            return EXIT_FAILURE;
        }
    }
    if (log_path) {
        log_fp = fopen(log_path, "a");
        if (!log_fp) { perror("fopen log"); return EXIT_FAILURE; }
    }

    signal(SIGINT,  handle_sigint);
    signal(SIGPIPE, SIG_IGN);

    server_fd = socket(AF_INET, SOCK_STREAM, 0);
    if (server_fd < 0) { perror("socket"); return EXIT_FAILURE; }

    int opt = 1;
    setsockopt(server_fd, SOL_SOCKET, SO_REUSEADDR, &opt, sizeof(opt));

    struct sockaddr_in addr = {
        .sin_family      = AF_INET,
        .sin_addr.s_addr = INADDR_ANY,
        .sin_port        = htons(port),
    };
    if (bind(server_fd, (struct sockaddr *)&addr, sizeof(addr)) < 0) {
        perror("bind"); return EXIT_FAILURE;
    }
    if (listen(server_fd, 16) < 0) {
        perror("listen"); return EXIT_FAILURE;
    }
    printf("Chat Hub listening on port %d...\n", port);

    for (;;) {
        struct sockaddr_in cli;
        socklen_t cl = sizeof(cli);
        int fd = accept(server_fd, (struct sockaddr *)&cli, &cl);
        if (fd < 0) {
            if (errno == EINTR) continue;
            perror("accept");
            break;
        }
        struct timeval to = { .tv_sec = 5, .tv_usec = 0 };
        setsockopt(fd, SOL_SOCKET, SO_SNDTIMEO, &to, sizeof(to));

        pthread_t tid;
        if (pthread_create(&tid, NULL, client_thread,
                           (void *)(intptr_t)fd) != 0) {
            perror("pthread_create");
            close(fd);
            continue;
        }
        pthread_detach(tid);
    }

    if (log_fp) fclose(log_fp);
    close(server_fd);
    return 0;
}
\end{lstlisting}

\subsection{\texttt{client.c}}

\begin{lstlisting}[caption={\texttt{client.c}, client interactif du hub}]
#include <stdio.h>
#include <stdlib.h>
#include <stdint.h>
#include <string.h>
#include <unistd.h>
#include <signal.h>
#include <pthread.h>
#include <netdb.h>
#include <arpa/inet.h>

#include "common.h"

static volatile int running = 1;

typedef struct { line_reader_t *lr; } recv_arg_t;

static void redraw_prompt(void) { printf("> "); fflush(stdout); }

static void *recv_thread(void *arg) {
    recv_arg_t *a = arg;
    char line[BUF_SIZE];

    while (running) {
        int r = lr_next(a->lr, line, sizeof(line));
        if (r <= 0) {
            if (running) {
                printf("\r\033[K[disconnected from server]\n");
                fflush(stdout);
                running = 0;
            }
            return NULL;
        }
        printf("\r\033[K");
        if (strncmp(line, "MSG ", 4) == 0) {
            const char *rest = line + 4;
            const char *sp   = strchr(rest, ' ');
            if (sp)
                printf("%.*s: %s\n", (int)(sp - rest), rest, sp + 1);
        } else if (strncmp(line, "PRIV ", 5) == 0) {
            const char *rest = line + 5;
            const char *sp   = strchr(rest, ' ');
            if (sp)
                printf("[private from %.*s] %s\n",
                       (int)(sp - rest), rest, sp + 1);
        } else if (strncmp(line, "JOIN ", 5) == 0) {
            printf("[*] %s a rejoint le chat\n", line + 5);
        } else if (strncmp(line, "LEAVE ", 6) == 0) {
            printf("[*] %s a quitte le chat\n", line + 6);
        } else if (strcmp(line, "LIST") == 0
                   || strncmp(line, "LIST ", 5) == 0) {
            printf("[users]%s\n", line + 4);
        } else if (strncmp(line, "ERR ", 4) == 0) {
            printf("[error] %s\n", line + 4);
        } else if (strcmp(line, "OK") != 0) {
            printf("%s\n", line);
        }
        redraw_prompt();
    }
    return NULL;
}

static int handshake(int fd, line_reader_t *lr,
                     char *out_name, size_t cap) {
    char nick[NAME_SIZE], line[BUF_SIZE], msg[BUF_SIZE];
    for (;;) {
        printf("Pseudo: ");
        fflush(stdout);
        if (!fgets(nick, sizeof(nick), stdin)) return -1;
        nick[strcspn(nick, "\n")] = '\0';
        if (nick[0] == '\0') continue;
        snprintf(msg, sizeof(msg), "NICK %s", nick);
        if (send_line(fd, msg) < 0) return -1;
        int r = lr_next(lr, line, sizeof(line));
        if (r <= 0) return -1;
        if (strcmp(line, "OK") == 0) {
            strncpy(out_name, nick, cap - 1);
            out_name[cap - 1] = '\0';
            return 0;
        }
        if (strncmp(line, "ERR ", 4) == 0)
            printf("[error] %s\n", line + 4);
    }
}

static int connect_to(const char *host, int port) {
    char portstr[16];
    snprintf(portstr, sizeof(portstr), "%d", port);
    struct addrinfo hints = { .ai_family   = AF_INET,
                              .ai_socktype = SOCK_STREAM };
    struct addrinfo *res;
    if (getaddrinfo(host, portstr, &hints, &res) != 0) return -1;
    int fd = socket(res->ai_family, res->ai_socktype, res->ai_protocol);
    if (fd < 0 || connect(fd, res->ai_addr, res->ai_addrlen) < 0) {
        if (fd >= 0) close(fd);
        freeaddrinfo(res);
        return -1;
    }
    freeaddrinfo(res);
    return fd;
}

int main(int argc, char *argv[]) {
    const char *host = NULL;
    int         port = 0;

    for (int i = 1; i < argc; i++) {
        if (strcmp(argv[i], "--server") == 0 && i + 1 < argc)
            host = argv[++i];
        else if (strcmp(argv[i], "--port") == 0 && i + 1 < argc)
            port = atoi(argv[++i]);
        else {
            fprintf(stderr, "Usage: %s --server <host> --port <port>\n",
                    argv[0]);
            return EXIT_FAILURE;
        }
    }
    if (!host || port <= 0 || port > 65535) {
        fprintf(stderr, "Usage: %s --server <host> --port <port>\n",
                argv[0]);
        return EXIT_FAILURE;
    }

    signal(SIGPIPE, SIG_IGN);

    int fd = connect_to(host, port);
    if (fd < 0) return EXIT_FAILURE;

    line_reader_t lr;
    lr_init(&lr, fd);

    char name[NAME_SIZE];
    if (handshake(fd, &lr, name, sizeof(name)) != 0) {
        close(fd); return EXIT_FAILURE;
    }
    printf("[INFO] Connecte en tant que %s\n", name);
    printf("Commandes: /msg <pseudo> <texte> | /users | /quit\n");
    redraw_prompt();

    pthread_t tid;
    recv_arg_t ra = { .lr = &lr };
    if (pthread_create(&tid, NULL, recv_thread, &ra) != 0) {
        perror("pthread_create"); close(fd); return EXIT_FAILURE;
    }

    char line[BUF_SIZE - NAME_SIZE - 16];
    while (running && fgets(line, sizeof(line), stdin)) {
        line[strcspn(line, "\n")] = '\0';
        if (line[0] == '\0') { redraw_prompt(); continue; }

        if (strcmp(line, "/quit") == 0) {
            send_line(fd, "QUIT"); running = 0; break;
        }
        if (strcmp(line, "/users") == 0) {
            send_line(fd, "LIST"); continue;
        }
        if (strncmp(line, "/msg ", 5) == 0) {
            const char *rest = line + 5;
            const char *sp   = strchr(rest, ' ');
            if (!sp || sp == rest) {
                printf("Usage: /msg <pseudo> <texte>\n");
                redraw_prompt(); continue;
            }
            char msg[BUF_SIZE];
            snprintf(msg, sizeof(msg), "PRIV %.*s %s",
                     (int)(sp - rest), rest, sp + 1);
            send_line(fd, msg); redraw_prompt(); continue;
        }
        if (line[0] == '/') {
            printf("Commande inconnue. /msg, /users, /quit\n");
            redraw_prompt(); continue;
        }
        char msg[BUF_SIZE];
        snprintf(msg, sizeof(msg), "MSG %s", line);
        send_line(fd, msg); redraw_prompt();
    }

    running = 0;
    shutdown(fd, SHUT_RDWR);
    pthread_join(tid, NULL);
    close(fd);
    return 0;
}
\end{lstlisting}

\subsection{Makefile}

\begin{lstlisting}[language=make, caption={Makefile, compilation du hub et du client}]
CC = gcc
CFLAGS = -Wall -Wextra -pthread

all: chat_serverd chat_client

chat_serverd: server.c common.h
	$(CC) $(CFLAGS) -o $@ server.c

chat_client: client.c common.h
	$(CC) $(CFLAGS) -o $@ client.c

clean:
	rm -f chat_serverd chat_client

.PHONY: all clean
\end{lstlisting}

% ── Comparaison des trois projets ─────────────────────────────────────────────
\section{Comparaison des trois projets}

\begin{center}
\begin{tabular}{@{}llll@{}}
\toprule
Aspect & Sockets 1 & Sockets 2 & Sockets 3 (Hub) \\
\midrule
Clients simultanés & 1 & 1 & Jusqu'à 64 \\
Mode de communication & Semi-duplex & Full-duplex & Full-duplex (salon) \\
Threads côté serveur & 0 & 1 par connexion & 1 par client (détaché) \\
Gestion des threads &, & \texttt{pthread\_join} & \texttt{pthread\_detach} \\
Protocole & Flux brut & Flux brut & Lignes texte structurées \\
Messages privés & Non & Non & Oui (\texttt{PRIV}) \\
Liste des connectés & Non & Non & Oui (\texttt{LIST}) \\
Journal persistant & Non & Non & Optionnel (\texttt{--log}) \\
Résolution DNS & Non & Non & Oui (\texttt{getaddrinfo}) \\
\bottomrule
\end{tabular}
\end{center}

% ── Bilan ─────────────────────────────────────────────────────────────────────
\section{Bilan}

Ce projet illustre les défis supplémentaires qu'introduit le passage d'un chat
à deux à un salon multi-clients : la nécessité de protéger un état partagé
(tableau des clients) par des mutex, l'importance de ne pas maintenir le verrou
pendant les envois réseau pour ne pas bloquer les autres fils, et l'utilisation
de \texttt{pthread\_detach} comme alternative à \texttt{pthread\_join} quand le
fil principal n'a pas besoin d'attendre la terminaison de chaque client.

L'introduction d'un protocole texte structuré est également une avancée
significative : elle découple clairement le format fil-à-fil de la logique
applicative et rend le serveur interopérable avec n'importe quel client
capable de parler le protocole, y compris \texttt{telnet} ou \texttt{netcat}.

\end{document}
